While \RLM{}s show strong performance on tasks beyond the context window limitations of existing LMs at reasonable inference costs, evaluations for more difficult and natural long-context processing tasks and the best mechanisms for implementing guardrails for \RLM{}s both remain highly under-explored. Broadly, \RLM{}s add a layer of complexity on top of existing LMs that may lead to unintentional side-effects like exploding sub-call costs, which we leave for future work to solve. We also note that future strategies involving asynchronous sub-calls and sandboxed REPLs can potentially significantly reduce the runtime and inference cost of \RLM{}s, but further contribute to this complexity. We include additional limitations and negative results in Appendix~\ref{appx0:didnt-work}. 


Lastly, we focused our experiments on evaluating \RLM{}s using \textit{existing} frontier models, but show initial evidence on a Qwen3-8B model that explicit training as a \RLM{} provides very rapid performance improvements, even outside the training domain. We hypothesize that \RLM{} trajectories can be viewed as a form of reasoning~\citep{openai2024openaio1card,deepseekai2025deepseekr1incentivizingreasoningcapability}, which can be trained by bootstrapping existing models~\citep{zelikman2022starbootstrappingreasoningreasoning, zelikman2024quietstarlanguagemodelsteach}. We hope that training native \RLM{}s can be treated as a new axis of scale to improve LM performance on general and long-horizon tasks.
